\documentclass[]{UCD_CS_MSc_AAI_Report}

% -----------------------------------------------------------------------------
% Packages
% -----------------------------------------------------------------------------
\usepackage{graphicx}
\usepackage{listings}
\usepackage{hyperref}
\usepackage{amsmath}
\usepackage{amssymb}
\usepackage{booktabs}
\usepackage{array}
\usepackage{xcolor}
\usepackage{braket}
\usepackage{subcaption}
\usepackage{rotating}

\usepackage[backend=biber]{biblatex}
\addbibresource{references.bib}

% -----------------------------------------------------------------------------
% Draft research-note commands
% Remove the coloured notes and planning lists before final submission.
% -----------------------------------------------------------------------------
\newcommand{\motive}[1]{{\color{blue!70!cyan}#1}}
\newcommand{\challenge}[1]{{\color{red}#1}}
\newcommand{\limitation}[1]{{\color{orange}#1}}
\newcommand{\method}[1]{{\color{cyan!80!black}#1}}
\newcommand{\result}[1]{{\color{green!60!black}#1}}
\newcommand{\benchmark}[1]{{\color{violet}#1}}
\newcommand{\comparison}[1]{{\color{pink!70!violet}#1}}
\newcommand{\theory}[1]{{\color{gray}#1}}
\newcommand{\tooling}[1]{{\color{brown}#1}}
\newcommand{\gap}[1]{{\color{black}#1}}
\newcommand{\validation}[1]{{\color{yellow!60!olive}#1}}

\newcommand{\sectionbrief}{\paragraph{Key points to cover.}}
\newcommand{\researchnotes}{\paragraph{Research notes and evidence.}}
\newcommand{\resultplaceholder}{\paragraph{Results to insert.}}

\hypersetup{
    colorlinks=true,
    linkcolor=blue,
    filecolor=magenta,
    urlcolor=cyan,
}

\lstset{
    mathescape=true,
    basicstyle=\ttfamily\small,
    breaklines=true,
    frame=single,
    numbers=left,
    numberstyle=\tiny,
    keywordstyle=\color{blue},
    commentstyle=\color{gray},
    language=Python
}

% -----------------------------------------------------------------------------
% Project details
% -----------------------------------------------------------------------------
\def\studentname{Edward Mallia}
\def\studentid{25233259}
\def\projecttitle{Quantum Ensemble Stacking with State Preservation}
\def\supervisorname{Dr Simon Caton}

\begin{document}
\maketitle

\tableofcontents
\newpage

% =============================================================================
% ABSTRACT
% =============================================================================
\abstract

% Write the abstract last. Cover each point in one or two sentences.
\begin{itemize}
    \item Context: practical limitations of large variational quantum classifiers and the motivation for ensembles of small circuits.
    \item Problem: conventional stacking measures each quantum base learner before aggregation, potentially creating an information bottleneck.
    \item Aim: design and evaluate a progression from measured ensemble aggregation to classical stacking, quantum re-encoding and measurement-free coherent stacking.
    \item Method: configurable VQC base learners, ECOC decomposition, systematic VQC architecture selection, controlled aggregation comparisons and multiclass tabular datasets.
    \item Evaluation: accuracy, macro-F1, class-level behaviour, convergence, training time and circuit-resource requirements.
    \item Principal findings: insert only after the final experiments and analysis are complete.
    \item Final conclusion: state whether coherent stacking was feasible, beneficial and practically justified relative to measured baselines.
\end{itemize}

% Optional front matter to add when needed:
% \chapter*{List of Abbreviations}
% QML, VQC, QNN, NISQ, ECOC, OOF, SVC, CZ, RZZ, etc.

% =============================================================================
% CHAPTER 1: INTRODUCTION
% =============================================================================
\chapter{Introduction}
\label{chap:introduction}

\section{Background and Motivation}
\label{sec:intro-background}

\sectionbrief
\begin{itemize}
    \item Introduce QML and variational quantum models as a near-term approach to supervised learning.
    \item Explain the practical restrictions of current quantum models: qubit limits, circuit depth, noise, expensive simulation and difficult optimisation.
    \item Motivate ensembles of several small VQCs as an alternative to one large monolithic circuit.
    \item Introduce stacking as a learned aggregation mechanism rather than fixed voting or averaging.
    \item Establish the specifically quantum issue: intermediate measurement converts a quantum state into classical outputs before the meta-learner can act on it.
    \item Motivate the central investigation without assuming that unmeasured states necessarily improve classification.
    \item End by narrowing the discussion from general QML to measurement-free quantum stacking.
\end{itemize}

\researchnotes
\begin{itemize}
    \motive{\item Ensemble methods combine complementary information from multiple learners and can improve accuracy, stability and generalisation on complex, noisy, imbalanced or high-dimensional data~\cite{dong_survey_2020}.}
    \motive{\item Stacking learns a second-level combiner from base-model predictions, making it a suitable framework for comparing classical and quantum aggregation mechanisms~\cite{demir_modified_2018}.}
    \motive{\item Quantum ensembles can retain the contributions of multiple classifiers rather than selecting one final model, motivating ensembles of simple quantum learners~\cite{schuld_quantum_2018}.}
    \motive{\item Small or shallow quantum learners have been proposed as a modular response to NISQ-era limits in circuit depth, trainability and noise~\cite{friedrich_quantum_2025,park_addressing_2025,incudini_resource_2023}.}
    \motive{\item Classical ensembles of small quantum classifiers have improved measured performance in several simulated and hardware-oriented studies, supporting ensemble construction as a relevant QML direction~\cite{song_quantum_2024,mcfarthing_classical_2024,tolotti_ensembles_2024}.}
    \challenge{\item QML remains limited by data loading, output extraction, benchmarking, noise and trainability, so the dissertation should avoid broad or unsupported quantum-advantage claims~\cite{biamonte_quantum_2017,du_quantum_2025}.}
\end{itemize}

\section{Problem Statement}
\label{sec:intro-problem}

\sectionbrief
\begin{itemize}
    \item Define conventional measured quantum stacking: independently measure each VQC, concatenate its classical outputs and train a meta-learner.
    \item Explain that a quantum meta-learner operating on re-encoded probabilities remains downstream of the same measurement boundary.
    \item Define the proposed alternative: preserve each trained base circuit as an unmeasured subcircuit and append a trainable meta-circuit before final measurement.
    \item State the uncertainty being tested: whether the preserved state contains exploitable information that is useful for classification.
    \item Identify the competing practical cost: wider and deeper circuits, more two-qubit gates, harder optimisation and increased simulation or hardware demands.
    \item Frame the dissertation as a controlled empirical comparison rather than a proof that measurement necessarily harms ensemble performance.
\end{itemize}

\researchnotes
\begin{itemize}
    \theory{\item An $n$-qubit state is described by $2^n$ amplitudes, while measurement yields classical outcomes and does not directly reveal the full state description~\cite{resch_quantum_2019}.}
    \theory{\item Entanglement stores information in relationships between qubits, motivating investigation of whether cross-register operations can exploit relationships unavailable after independent base measurements~\cite{resch_quantum_2019}.}
    \challenge{\item Quantum states are fragile and cannot be freely copied or persistently stored, making quantum intermediate representations fundamentally different from saved classical stacking features~\cite{resch_quantum_2019}.}
    \challenge{\item Increased circuit size and global entanglement can worsen optimisation, noise exposure and barren-plateau behaviour~\cite{du_quantum_2025,cerezo_challenges_2022}.}
    \gap{\item Existing work does not provide a direct controlled comparison of classical measured stacking, quantum re-encoding and a coherent meta-circuit applied to trained VQC base states.}
\end{itemize}

\section{Aim}
\label{sec:intro-aim}

\sectionbrief
\begin{itemize}
    \item State one overall aim: to design, implement and evaluate a quantum ensemble architecture whose aggregation stage progresses from measured methods to a fully quantum, measurement-free stacking condition.
    \item Emphasise that the base VQCs and ECOC decomposition provide a controlled foundation; the primary research focus is the stacking boundary and meta-level aggregation.
\end{itemize}

\section{Research Objectives}
\label{sec:intro-objectives}

\begin{enumerate}
    \item Develop configurable and trainable VQC base learners for small multiclass tabular datasets.
    \item Investigate VQC architecture choices and select a compact set of feasible base-learner designs.
    \item Construct a controlled ECOC ensemble of simple quantum learners.
    \item Establish direct measured ECOC aggregation as the basic quantum ensemble baseline.
    \item Implement classical measured stacking using stratified out-of-fold meta-training.
    \item Implement a quantum meta-learner operating on re-encoded measured base outputs.
    \item Develop a coherent architecture that removes intermediate measurements and appends a trainable quantum meta-circuit to frozen base circuits.
    \item Compare predictive performance, optimisation behaviour, computational cost and circuit-resource requirements across aggregation conditions.
    \item Analyse which dataset and ensemble characteristics are associated with any observed stacking benefit.
\end{enumerate}

\section{Research Questions}
\label{sec:intro-rqs}

\begin{enumerate}
    \item Can an ensemble of small VQC base learners provide a viable foundation for multiclass quantum classification?
    \item How does classical stacking of measured VQC outputs compare with direct measured ECOC aggregation?
    \item Does a quantum meta-learner trained on re-encoded measured outputs provide an advantage over a classical meta-learner?
    \item Can trained VQC base learners be composed into a trainable measurement-free quantum stacking circuit?
    \item Does preserving the base learners' quantum states improve classification performance, macro-F1, robustness or convergence relative to measured stacking?
    \item What additional computational and quantum-resource costs are introduced by measurement-free stacking?
    \item For which dataset or ensemble characteristics, if any, is coherent quantum stacking most beneficial?
\end{enumerate}

\section{Expected and Actual Contributions}
\label{sec:intro-contributions}

\sectionbrief
\begin{itemize}
    \item A configurable framework for VQC architecture construction and evaluation.
    \item A systematic VQC architecture-search pipeline and cross-dataset ranking process.
    \item A controlled ECOC framework for generating simple multiclass quantum base learners.
    \item A strong out-of-fold classical stacking baseline for measured VQC ensembles.
    \item A quantum re-encoding comparison that isolates the effect of the meta-learner type.
    \item A measurement-free architecture formed by composing frozen VQC base circuits with a trainable coherent meta-circuit.
    \item An empirical analysis of measurement as a possible base-to-meta information bottleneck.
    \item A performance--cost analysis covering prediction quality, training cost and circuit resources.
    \item Rewrite this list near submission so that it states only contributions that were actually completed and evaluated.
\end{itemize}

\section{Scope and Delimitations}
\label{sec:intro-scope}

\sectionbrief
\begin{itemize}
    \item The central contribution concerns aggregation, not a novel ECOC coding theory or a general-purpose new VQC.
    \item The study focuses on supervised multiclass classification using small and moderate tabular datasets.
    \item ECOC is used as practical scaffolding for repeatable low-complexity base tasks.
    \item Base learners are deliberately small so that the stacking mechanism remains the main experimental variable.
    \item Most experiments are expected to rely on simulation because coherent statevector cost scales exponentially with total qubit count.
    \item Include noisy simulation or IBM hardware only if those experiments are actually completed.
    \item Do not claim quantum speed-up or quantum advantage unless directly supported by the final evidence.
\end{itemize}

\section{Dissertation Structure}
\label{sec:intro-structure}

\sectionbrief
\begin{itemize}
    \item Summarise Chapters~\ref{chap:related-work}--\ref{chap:conclusion} in one short paragraph after the chapters have stabilised.
\end{itemize}

% =============================================================================
% CHAPTER 2: RELATED WORK
% =============================================================================
\chapter{Related Work}
\label{chap:related-work}

\section{Ensemble Learning}
\label{sec:rw-ensembles}

\subsection{Motivation for Ensemble Learning}

\sectionbrief
\begin{itemize}
    \item Explain why multiple imperfect learners can outperform or stabilise a single learner.
    \item Introduce the accuracy--diversity trade-off.
    \item Identify sources of diversity: data samples, feature subsets, model architectures, class decompositions, initialisation and prediction behaviour.
    \item Relate ensemble diversity to the dissertation's feature subspacing, bagging, ECOC tasks and VQC-template variation.
\end{itemize}

\researchnotes
\begin{itemize}
    \method{\item Ensemble diversity can be introduced through sample subsets, feature subspaces, classifier selection, weighting and ensemble-level optimisation~\cite{dong_survey_2020}.}
    \comparison{\item Ensemble quality depends on the joint behaviour of accuracy, diversity, generalisation, stability, sparsity and complexity rather than constituent accuracy alone~\cite{dong_survey_2020}.}
    \method{\item Random-subspace methods provide direct precedent for assigning different feature subsets to VQC base learners~\cite{dong_survey_2020,radojicic_reducing_2026}.}
    \result{\item Prior VQC composition work identifies prediction disagreement as a stronger indicator of ensemble contribution than individual F1 or architectural difference, supporting analysis of actual base-prediction diversity~\cite{pinto_composing_2026}.}
    \method{\item Diversity-aware QBoost work explicitly optimises disagreement alongside prediction loss and regularisation, providing precedent for treating diversity as a design objective rather than only a post-hoc statistic~\cite{liu_construction_2025}.}
\end{itemize}

\subsection{Bagging, Boosting and Stacking}

\sectionbrief
\begin{itemize}
    \item Distinguish parallel bagging, sequential boosting and two-level stacking.
    \item Explain why stacking is the principal aggregation method in this dissertation.
    \item Describe fixed aggregation, weighted aggregation and learned aggregation.
    \item Clarify that bagging is also used within the project as a source of base-learner diversity rather than as the main research contribution.
\end{itemize}

\researchnotes
\begin{itemize}
    \method{\item Classical supervised ensembles commonly include bagging, boosting, random forests, random subspaces and stacking~\cite{dong_survey_2020}.}
    \method{\item Stacking uses base predictions as second-level input features and trains a combiner to learn how those predictions should be integrated~\cite{demir_modified_2018}.}
    \comparison{\item Different stackers can favour different class-level behaviours, showing that the meta-learner is a substantive design choice rather than neutral plumbing~\cite{demir_modified_2018}.}
    \result{\item Prior VQC ensemble work reports stronger gains from stacking than from bagging, supporting the dissertation's emphasis on learned aggregation~\cite{pinto_composing_2026}.}
\end{itemize}

\subsection{Out-of-Fold Stacking}

\sectionbrief
\begin{itemize}
    \item Explain why training the meta-learner on in-sample base predictions creates leakage and overly optimistic performance.
    \item Define stratified $K$-fold out-of-fold prediction generation.
    \item Explain the distinction between temporary fold ensembles and the final base ensemble trained on the full training partition.
    \item Connect this directly to the implemented \texttt{StackedECOC} training procedure.
\end{itemize}

\researchnotes
\begin{itemize}
    \method{\item Five-fold stacking predictions are used in prior hybrid ensemble work to create valid learned features before downstream quantum classification~\cite{xiong_hybrid_2024}.}
    \method{\item Modified stacking uses feature-subspace base models and validation-driven model selection before training the combiner, providing precedent for separating base-model generation from meta-training~\cite{demir_modified_2018}.}
\end{itemize}

\subsection{Error-Correcting Output Codes}

\sectionbrief
\begin{itemize}
    \item Introduce ECOC as a method for decomposing multiclass classification into several simpler learner tasks.
    \item Define the coding matrix, class codewords, learner columns or rows according to the implementation convention, and decoding.
    \item Explain binary ECOC and the optional generalisation to greater decomposition depth.
    \item Explain why ECOC is appropriate here: controllable learner count, simple VQC tasks and consistent use across datasets.
    \item State explicitly that ECOC is not the central novelty.
\end{itemize}

\researchnotes
\begin{itemize}
    \motive{\item ECOC decomposes a $c$-class problem into several simpler learner problems and decodes the resulting prediction codeword against class codewords~\cite{nguyen_optimal_2021}.}
    \theory{\item Strong ECOC matrices aim to maintain separation between class codewords while avoiding redundant or constant learners~\cite{nguyen_optimal_2021}.}
    \method{\item $N$-ary ECOC generalises binary class partitions but increases the complexity of each base task, supporting binary decomposition as the principal dissertation setting~\cite{nguyen_optimal_2021}.}
    \gap{\item Classical ECOC optimisation literature does not address VQC trainability, qubit budgets or coherent quantum stacking~\cite{nguyen_optimal_2021}.}
\end{itemize}

\section{Quantum Machine Learning}
\label{sec:rw-qml}

\subsection{Near-Term Quantum Machine Learning}

\sectionbrief
\begin{itemize}
    \item Distinguish fault-tolerant quantum algorithms from NISQ-compatible hybrid variational approaches.
    \item Introduce the hybrid quantum--classical optimisation loop.
    \item Explain why VQCs are appropriate to the dissertation's implementation and resource constraints.
    \item Treat quantum computing as a specialised accelerator or experimental computational model rather than a wholesale replacement for classical ML.
\end{itemize}

\researchnotes
\begin{itemize}
    \comparison{\item Fault-tolerant QML may support stronger theoretical speed-ups but requires error-corrected devices, whereas QNNs and kernels are more experimentally accessible in the NISQ setting~\cite{du_quantum_2025}.}
    \method{\item Near-term variational algorithms repeatedly execute a parameterised quantum circuit, measure its outputs and update parameters through a classical optimiser~\cite{resch_quantum_2019,du_quantum_2025}.}
    \limitation{\item Input preparation, output reconstruction and practical resource estimates can undermine theoretical QML speed-up claims~\cite{biamonte_quantum_2017}.}
    \comparison{\item Early QML literature distinguishes quantum acceleration of classical subroutines from genuinely quantum learning models, supporting the dissertation's focus on the learning and aggregation mechanism itself~\cite{schuld_introduction_2015}.}
\end{itemize}

\subsection{Variational Quantum Classifiers}

\sectionbrief
\begin{itemize}
    \item Define a VQC as feature encoding, parameterised quantum transformations, entangling operations, measurement and classical optimisation.
    \item Describe how measurements are converted into class probabilities in this project.
    \item Introduce local versus fuller measurement choices.
    \item Explain why architecture selection matters for expressivity, trainability and generalisation.
\end{itemize}

\researchnotes
\begin{itemize}
    \method{\item A standard QNN workflow consists of encoding, trainable quantum processing, measurement, loss evaluation and classical parameter updates~\cite{du_quantum_2025}.}
    \theory{\item Circuit design is governed by the interaction between expressivity, trainability and generalisation~\cite{du_quantum_2025}.}
    \challenge{\item VQCs can be sensitive to initialisation, dataset size, circuit depth and cost-function concentration~\cite{vasan_exploring_2025,friedrich_quantum_2025}.}
\end{itemize}

\subsection{Quantum Feature Encoding and Re-uploading}

\sectionbrief
\begin{itemize}
    \item Introduce angle encoding and the mapping of classical features to rotation gates.
    \item Explain feature density, feature assignment across qubits and repeated data re-uploading.
    \item Discuss the trade-off between richer feature representation and increased depth or gate count.
    \item Introduce pairwise feature interactions through RZZ-based encoding where relevant.
\end{itemize}

\researchnotes
\begin{itemize}
    \method{\item Data re-uploading is a recognised strategy for increasing the expressivity of classical-data quantum feature maps~\cite{du_quantum_2025}.}
    \method{\item Repeated feature-encoding rotations can increase accessible Fourier frequencies and classifier expressivity~\cite{mcfarthing_classical_2024}.}
    \comparison{\item Feature-map ordering and the placement of trainable layers can materially change trainability and eliminate redundant gates~\cite{salmenpera_impact_2024}.}
    \challenge{\item Higher feature dimensionality increases circuit width or depth and can make full-feature VQC training prohibitive, motivating feature subspacing~\cite{radojicic_reducing_2026}.}
\end{itemize}

\subsection{Ansatz Design and Entanglement}

\sectionbrief
\begin{itemize}
    \item Describe parameterised rotation layers and entangling layers.
    \item Discuss CZ and RZZ entanglers, linear/ring/odd--even connectivity and hardware-aware design.
    \item Explain why the dissertation searches several qualitatively different shallow ansatz structures rather than committing to one arbitrary template.
    \item Relate the selected gate family to IBM Heron-compatible gates where appropriate.
\end{itemize}

\researchnotes
\begin{itemize}
    \challenge{\item Similar-depth ansatz families can differ substantially in performance, creating an empirical ``ansatz lottery'' if only one architecture is tested~\cite{pinto_composing_2026}.}
    \method{\item Shallow circuit families, restricted entanglement and hardware-efficient templates can improve practical trainability and reduce noise exposure~\cite{friedrich_quantum_2025,park_addressing_2025}.}
    \challenge{\item Limited hardware connectivity may introduce SWAP operations, increasing depth, latency and error exposure~\cite{resch_quantum_2019}.}
    \result{\item Several studies observe diminishing performance returns after a small number of trainable layers, supporting shallow initial designs~\cite{salmenpera_impact_2024,incudini_resource_2023}.}
\end{itemize}

\subsection{Trainability and Scaling}

\sectionbrief
\begin{itemize}
    \item Introduce barren plateaus, vanishing gradients, local minima and initialisation sensitivity.
    \item Explain how circuit width, depth, global measurements and entanglement can affect optimisation.
    \item Motivate shallow base learners, early stopping and staged training.
    \item Explain the exponential classical statevector cost of increasing the total coherent qubit count.
\end{itemize}

\researchnotes
\begin{itemize}
    \challenge{\item Barren plateaus can cause gradient magnitudes to decrease exponentially with system size~\cite{du_quantum_2025,cerezo_challenges_2022}.}
    \method{\item Initialisation, local cost functions, structured circuits and regularisation are common mitigation strategies~\cite{du_quantum_2025}.}
    \method{\item Partitioning one deep quantum layer into multiple shallow components can increase gradient variance and improve convergence in some settings~\cite{friedrich_quantum_2025}.}
    \challenge{\item More global entanglement may worsen trainability and generalisation, while local or modular entanglement can reduce optimisation difficulty~\cite{park_addressing_2025}.}
\end{itemize}

\subsection{Measurement and Quantum Information}

\sectionbrief
\begin{itemize}
    \item Explain measurement as the final conversion from quantum state to classical outcome probabilities.
    \item Distinguish the mathematical state from the limited statistics obtained through measurement.
    \item Explain why independently measuring base registers prevents later coherent cross-register operations.
    \item Avoid claiming that all discarded state information is useful for the target task.
    \item Establish measurement as the empirical variable examined at the base-to-meta boundary.
\end{itemize}

\researchnotes
\begin{itemize}
    \theory{\item Amplitudes are not directly observable and expectation values require repeated measurement or exact simulation~\cite{resch_quantum_2019,tolotti_ensembles_2024}.}
    \challenge{\item Measured ensemble pipelines are sensitive to shot noise and readout precision~\cite{tolotti_ensembles_2024,tolotti_weighted_2025}.}
    \theory{\item The no-cloning theorem and the absence of practical persistent quantum memory prevent treating quantum intermediate states like ordinary stored classical features~\cite{resch_quantum_2019}.}
\end{itemize}

\section{Quantum Ensemble Learning}
\label{sec:rw-quantum-ensembles}

\subsection{Coherent Superposition Ensembles}

\sectionbrief
\begin{itemize}
    \item Review early definitions of quantum ensembles based on superpositions over classifier parameters or indexed internal classifiers.
    \item Explain the theoretical appeal of parallel evaluation and amplitude-weighted aggregation.
    \item Distinguish these methods from the dissertation's finite set of independently trained VQC subcircuits.
\end{itemize}

\researchnotes
\begin{itemize}
    \method{\item Schuld and Petruccione prepare a weighted superposition over classifier parameter states and extract a collective decision through final measurement~\cite{schuld_quantum_2018}.}
    \method{\item Araujo and da Silva compare untrained and trained classifier ensembles and find that optimisation is important on benchmark datasets~\cite{araujo_quantum_2020}.}
    \method{\item Weighted indexed ensembles can generate different internal data compositions in superposition and re-encode classically learned aggregation weights for coherent test-time execution~\cite{tolotti_weighted_2025}.}
    \limitation{\item These approaches rely on idealised state preparation, indexed data access, post-selection or simulated amplitudes, limiting direct equivalence to near-term trained VQC stacks~\cite{schuld_quantum_2018,araujo_quantum_2020,tolotti_weighted_2025}.}
\end{itemize}

\subsection{Measured Quantum Bagging and Boosting}

\sectionbrief
\begin{itemize}
    \item Review measured quantum base learners combined through voting, averaging, bagging or boosting.
    \item Explain what these approaches demonstrate about ensemble robustness and small-circuit modularity.
    \item Identify that they still measure base learners before aggregation.
\end{itemize}

\researchnotes
\begin{itemize}
    \method{\item AdaBoost.Q uses measured class probabilities as confidence information for reweighting samples and combining quantum classifiers~\cite{song_quantum_2024}.}
    \method{\item Bagging and boosting of compact QNN or single-qubit classifiers can improve predictive performance while reducing matched training or circuit costs~\cite{mcfarthing_classical_2024,incudini_resource_2023}.}
    \method{\item Bagging--QSVC and related measured quantum ensembles aggregate independently measured predictions classically~\cite{abdulsalam_explainable_2023}.}
    \method{\item Quantum bootstrapped bagging combines quantum-inspired subsampling or clustering with measured or classical aggregation, providing robustness evidence but not coherent VQC meta-learning~\cite{rathi_quantum_2025}.}
    \comparison{\item These methods provide evidence that ensemble construction can help quantum classifiers, but they do not preserve base states for a coherent meta-circuit.}
\end{itemize}

\subsection{Measured and Hybrid Quantum Stacking}

\sectionbrief
\begin{itemize}
    \item Review stacking systems that use quantum learners as measured base models and a classical stacker.
    \item Review systems that combine classical and quantum learners in a classical stacking architecture.
    \item Identify their relevance as measured baselines rather than coherent architectures.
\end{itemize}

\researchnotes
\begin{itemize}
    \method{\item Vasan et al. combine measured QSVC, QVC and QNN posterior probabilities with a Random Forest meta-learner~\cite{vasan_exploring_2025}.}
    \method{\item Kang and Shin combine measured ML, DNN and QNN predictions through a classical meta-model~\cite{kang_amino_2025}.}
    \method{\item Xiong et al. first use classical stacking for feature reduction and then apply a heterogeneous measured quantum voting ensemble~\cite{xiong_hybrid_2024}.}
    \result{\item These studies report improved aggregate performance or minority-class behaviour, supporting measured stacking as a substantive baseline~\cite{vasan_exploring_2025,kang_amino_2025,xiong_hybrid_2024}.}
    \gap{\item None trains a coherent quantum meta-circuit directly on the unmeasured states of trained VQC base learners.}
\end{itemize}

\subsection{Quantum Meta-Learning after Measurement}

\sectionbrief
\begin{itemize}
    \item Review systems that re-encode classical base predictions into a quantum kernel or variational meta-model.
    \item Explain that they test the value of a quantum meta-model but not preservation of the original base states.
    \item Use this distinction to motivate the dissertation's re-encoding comparison condition.
\end{itemize}

\researchnotes
\begin{itemize}
    \method{\item SEQSVM concatenates classical ensemble predictions into a compact meta-feature vector and feeds it to a quantum-kernel SVM~\cite{akrom_bloodbrain_2025}.}
    \method{\item Quantum stacking studies can provide the quantum meta-learner with predicted classes, confidence values or reduced high-level features~\cite{tolotti_ensembles_2024,akrom_bloodbrain_2025}.}
    \comparison{\item Re-encoding can introduce quantum processing at the meta-level, but cannot reconstruct information removed by the preceding base measurement.}
\end{itemize}

\subsection{Modular and Coherent Quantum Aggregation}

\sectionbrief
\begin{itemize}
    \item Review shallow ensemble layers, multi-chip VQCs, controlled-trajectory ensembles and amplitude-amplification aggregation.
    \item Separate approaches that preserve coherence from those that use modular circuits but aggregate classically.
    \item Explain where the proposed frozen-base coherent stack is similar and where it differs.
\end{itemize}

\researchnotes
\begin{itemize}
    \method{\item Friedrich and Maziero replace a deeper quantum layer with several independently measured shallow circuits whose outputs are summed~\cite{friedrich_quantum_2025}.}
    \method{\item Multi-chip ensemble VQCs partition inputs across small circuits but intentionally avoid cross-chip entanglement and combine measured outputs through a classical function~\cite{park_addressing_2025}.}
    \method{\item Rathi and Kumar create multiple transformed trajectories in superposition before a shared classifier and final ensemble measurement~\cite{rathi_enhanced_2025}.}
    \method{\item Safina uses amplitude amplification to estimate aggregate class probability from classically prepared classifier probabilities~\cite{safina_quantum_2023}.}
    \gap{\item These methods do not implement the same staged process of independently training VQC base learners, freezing their parameters and appending a trainable cross-register meta-VQC.}
\end{itemize}

\section{Gap Analysis}
\label{sec:rw-gap}

\subsection{Limitations of Existing Work}

\sectionbrief
\begin{itemize}
    \item Most practical quantum ensembles measure each component before aggregation.
    \item Hybrid and re-encoded quantum stackers operate only on classical outputs produced after collapse.
    \item Coherent ensemble proposals often rely on fault-tolerant assumptions, indexed data access, idealised state preparation or non-VQC classifier models.
    \item Modular VQC studies often avoid cross-register coupling for scalability and trainability.
    \item Existing studies rarely hold the base ensemble fixed while comparing direct measured aggregation, classical stacking, quantum re-encoding and coherent stacking.
    \item Many evaluations use small binary datasets or limited simulator settings and do not report both predictive and circuit-resource costs.
\end{itemize}

\researchnotes
\begin{itemize}
    \gap{\item Early quantum ensemble work explicitly leaves open alternative coherent definitions of quantum ensembles~\cite{schuld_quantum_2018}.}
    \gap{\item Measured quantum stacking studies do not investigate direct aggregation of unmeasured trained VQC states~\cite{vasan_exploring_2025,kang_amino_2025}.}
    \gap{\item Quantum meta-learning after re-encoding leaves the measurement boundary intact~\cite{akrom_bloodbrain_2025,tolotti_ensembles_2024}.}
    \gap{\item Shallow or multi-chip quantum ensembles improve modularity but still sum or classically combine measured outputs~\cite{friedrich_quantum_2025,park_addressing_2025}.}
    \gap{\item Quantum bagging frameworks may assume QRAM-style access or retain classical clustering and aggregation steps, leaving the trained coherent VQC stacking problem unresolved~\cite{rathi_quantum_2025}.}
    \gap{\item Supervisor-provided VQC ensemble studies identify VQC meta-learning as a natural extension but do not test coherent measurement-free stacking~\cite{radojicic_reducing_2026,pinto_composing_2026}.}
\end{itemize}

\subsection{Research Gap Addressed by This Dissertation}

\sectionbrief
\begin{itemize}
    \item State the gap in one precise paragraph: limited empirical evidence exists for stacking trained VQC base learners by preserving their unmeasured states and applying a trainable quantum meta-circuit across their registers.
    \item Explain that the study responds through a controlled sequence of aggregation conditions sharing the same basic ensemble foundation.
    \item Emphasise that the comparison evaluates both prediction quality and the additional cost of state preservation.
\end{itemize}

% =============================================================================
% CHAPTER 3: APPROACH AND METHODOLOGY
% =============================================================================
\chapter{Approach and Methodology}
\label{chap:methodology}

\section{Research Design}
\label{sec:method-research-design}

\sectionbrief
\begin{itemize}
    \item Describe the dissertation as a controlled empirical comparison of increasingly quantum aggregation mechanisms.
    \item Define the principal independent variable as aggregation strategy.
    \item Define dependent measures: predictive metrics, optimisation behaviour, runtime and quantum resources.
    \item Explain which factors are held consistent across conditions: datasets, preprocessing, outer splits, base-learner pool and evaluation procedure.
    \item Explain where exact equivalence is impossible, such as classical versus quantum meta-model capacity.
\end{itemize}

\subsection{Experimental Conditions}

\begin{enumerate}
    \item Classical reference models: SVC and Random Forest.
    \item Selected individual VQCs.
    \item Direct measured QuantumECOC aggregation.
    \item Classical measured StackedECOC using out-of-fold base outputs.
    \item Quantum meta-learning through re-encoded measured outputs.
    \item Measurement-free coherent stacking using frozen base circuits and a trainable quantum meta-circuit.
\end{enumerate}

\researchnotes
\begin{itemize}
    \comparison{\item The comparison should isolate the effect of aggregation rather than allowing every condition to use unrelated base learners or preprocessing.}
    \method{\item Prior ensemble design studies support varying feature subsets, learner selection and ensemble size in a bounded structured search rather than an unrestricted sweep~\cite{dong_survey_2020,demir_modified_2018,pinto_composing_2026}.}
\end{itemize}

\section{System Overview}
\label{sec:method-system}

\sectionbrief
\begin{itemize}
    \item Add an architecture figure showing: dataset $\rightarrow$ preprocessing $\rightarrow$ ECOC decomposition $\rightarrow$ VQC base learners $\rightarrow$ selected aggregation condition $\rightarrow$ final prediction.
    \item Separate the system into a base-ensemble level and a stacking/meta-learning level.
    \item Add a second figure contrasting the information flow in measured, re-encoded and coherent stacking.
    \item Explain the incremental implementation path from direct ensemble aggregation to the final coherent model.
\end{itemize}

\section{Implementation and Computational Environment}
\label{sec:method-environment}

\sectionbrief
\begin{itemize}
    \item Python version and core libraries: PyTorch, PennyLane, NumPy, pandas and scikit-learn.
    \item Custom VQC layer construction, forward pass, training/validation loop, prediction functions and result collection.
    \item Configurable ansatz definitions and template serialisation.
    \item Simulation devices and differentiation methods actually used.
    \item UCD SONIC HPC environment, Slurm execution and multiprocessing where relevant.
    \item Structured JSONL result logging, deterministic job identifiers and completed-job skipping.
    \item GitHub repository and reproducibility information.
    \item Include Qiskit or IBM Runtime only for work that is actually used in the final experiments.
\end{itemize}

\researchnotes
\begin{itemize}
    \tooling{\item PennyLane, PyTorch and Qiskit are established tools for hybrid quantum-classical experimentation~\cite{du_quantum_2025,resch_quantum_2019}.}
    \limitation{\item Simulator results represent best-case or modelled behaviour unless an explicit device-noise model or hardware execution is included~\cite{tolotti_ensembles_2024}.}
\end{itemize}

\section{Datasets}
\label{sec:method-datasets}

\subsection{Dataset Roles}

\sectionbrief
\begin{itemize}
    \item Define sanity-check datasets used to validate implementation and basic trainability.
    \item Define main evaluation datasets used for systematic comparisons.
    \item Define final stress-test datasets used on a reduced set of selected configurations.
    \item State the final list only after the full experiment plan is frozen.
\end{itemize}

\subsection{Dataset Characteristics}

% Replace placeholders with the final dataset table.
\begin{table}[ht]
    \centering
    \caption{Dataset characteristics and experimental roles.}
    \label{tab:datasets}
    \begin{tabular}{lrrrrll}
        \toprule
        Dataset & Samples & Features & Classes & Imbalance & Feature types & Role \\
        \midrule
        % Iris & ... & ... & ... & ... & ... & Sanity check \\
        % ...
        \bottomrule
    \end{tabular}
\end{table}

\sectionbrief
\begin{itemize}
    \item Discuss sample size, feature dimensionality, number of classes, imbalance, missing values and mixed data types.
    \item Explain why the set supports analysis across different levels of classification difficulty.
    \item Record any original unconventional train/test splits and justify whether they were retained or replaced.
\end{itemize}

\section{Data Preprocessing}
\label{sec:method-preprocessing}

\subsection{Cleaning and Missing Values}
\begin{itemize}
    \item Dataset-specific invalid-value handling.
    \item Numerical imputation strategy and fitted values.
    \item Categorical missing-value handling.
\end{itemize}

\subsection{Categorical Encoding}
\begin{itemize}
    \item Encoding strategy for nominal and ordinal variables.
    \item Avoidance of train--test leakage.
    \item Preservation of final feature mappings for reproducibility.
\end{itemize}

\subsection{Scaling}
\begin{itemize}
    \item Scaling range required by angle encoding, including the implemented $[0,\pi]$ range where applicable.
    \item Fit the scaler on training data and apply it to validation/test data.
    \item Explain any condition-specific scaling required by classical or quantum models.
\end{itemize}

\subsection{Data Splits}
\begin{itemize}
    \item Outer train/test split and stratification.
    \item Validation data used for early stopping or base-learner weighting.
    \item Inner out-of-fold splits used only for meta-feature construction.
    \item Repeated seeds or runs, where computationally feasible.
\end{itemize}

\section{Variational Quantum Classifier}
\label{sec:method-vqc}

\subsection{VQC Interface and Output Mapping}
\sectionbrief
\begin{itemize}
    \item Inputs: qubit count, class count, feature subset, ansatz template, measurement mode, device and differentiation method.
    \item Materialisation of the PennyLane quantum layer and integration with PyTorch.
    \item Conversion of measured basis-state probabilities into class probabilities.
    \item Minimum versus full measurement mode.
    \item Training, prediction, scoring and result-extraction functions.
\end{itemize}

\subsection{Feature Subspacing}
\begin{itemize}
    \item Selection of a subset of original features for each learner.
    \item Feature density as either a fraction or fixed number of features.
    \item Deterministic learner-specific seeds.
    \item Role in reducing circuit cost and increasing ensemble diversity.
\end{itemize}

\researchnotes
\begin{itemize}
    \result{\item Random-subspace VQC studies report that smaller feature ratios can preserve performance while reducing training cost, although the optimum is dataset-dependent~\cite{radojicic_reducing_2026}.}
    \benchmark{\item Resource-saving QNN ensembles can reduce qubit requirements substantially, but lower feature coverage may also reduce predictive information~\cite{incudini_resource_2023}.}
\end{itemize}

\subsection{Encoding Layers}
\begin{itemize}
    \item Angle-encoding layer and mapping of features to qubits.
    \item Feature strategies: same, shifted or explicit assignments.
    \item Feature re-uploading and feature count per qubit.
    \item Pairwise RZZ feature encoding where included.
    \item Barriers used only for circuit visualisation or transpilation clarity, if relevant.
\end{itemize}

\subsection{Trainable and Entangling Layers}
\begin{itemize}
    \item Parameterised RX and RZ rotations.
    \item CZ and trainable RZZ entanglers.
    \item Linear, ring, cyclic, odd--even or other implemented connectivity patterns.
    \item Placement of entanglement relative to encoding and trainable layers.
    \item Parameter count and gate-count derivation.
\end{itemize}

\section{VQC Architecture Exploration}
\label{sec:method-vqc-search}

\subsection{Purpose}
\begin{itemize}
    \item Explain why the final ensemble cannot be evaluated fairly using one arbitrarily selected circuit.
    \item Identify the target: compact, trainable and sufficiently expressive base architectures rather than maximum standalone accuracy at any resource cost.
\end{itemize}

\subsection{Search Space}
\begin{itemize}
    \item Qubit counts evaluated.
    \item Feature density and features per qubit.
    \item Minimum versus full measurement.
    \item Encoding style, feature strategy and re-upload count.
    \item Number and arrangement of trainable layers.
    \item Entangling-layer positions, connectivity and entangler type.
    \item Final exact values should be inserted from the executed job generator rather than memory.
\end{itemize}

\subsection{Search Execution}
\begin{itemize}
    \item Job-list generation and stable job identifiers.
    \item Parallel execution on local or HPC resources.
    \item Early stopping, epoch budget and validation protocol.
    \item JSONL logging and later reduction of redundant output fields.
    \item Failed-run handling and result repair where needed.
\end{itemize}

\subsection{Ranking and Selection}
\begin{itemize}
    \item Per-dataset ranking.
    \item Average rank or percentile across datasets.
    \item Overall configuration-level and ansatz-level aggregation.
    \item Performance consistency and resource cost.
    \item Shortlisting of qualitatively different high-performing base designs.
    \item Avoid selection using the final held-out test results.
\end{itemize}

\researchnotes
\begin{itemize}
    \challenge{\item Ansatz performance can vary strongly even for circuits with similar depth and parameter count, justifying systematic architecture selection~\cite{pinto_composing_2026}.}
    \method{\item Model-pool reuse separates expensive VQC training from later ensemble construction and analysis~\cite{pinto_composing_2026}.}
    \result{\item Shallow architectures can provide strong F1-per-second behaviour and reduce noise exposure~\cite{pinto_composing_2026}.}
\end{itemize}

\section{ECOC Base Ensemble}
\label{sec:method-ecoc}

\subsection{Coding Matrix Construction}
\begin{itemize}
    \item Define the implemented construction based on class count, learner count and ECOC depth.
    \item Explain how modular arithmetic generates learner-specific class groupings.
    \item Explain row/column orientation exactly as used in the implementation.
    \item Explain repetition or rolling when the requested learner count exceeds the initial matrix size.
\end{itemize}

\subsection{Learner Initialisation}
\begin{itemize}
    \item One VQC per ECOC subtask.
    \item Learner-specific class count based on its code partition.
    \item Template assignment and feature-subset selection.
    \item Controlled use of same, few-different and all-different base architectures.
\end{itemize}

\subsection{Training and Diversity}
\begin{itemize}
    \item Mapping original labels to each learner's ECOC targets.
    \item Stratified bagging or full-data training.
    \item Class weighting where used.
    \item Diversity from ECOC tasks, feature subsets, data samples, ansatz templates, entangler type and feature assignment.
\end{itemize}

\subsection{Direct Measured Aggregation}
\begin{itemize}
    \item Independently measure each VQC.
    \item Convert learner probabilities to scores over original classes using the ECOC matrix.
    \item Describe weighted or unweighted aggregation.
    \item Decode the highest final class score.
    \item Establish this as the direct measured quantum ensemble baseline.
\end{itemize}

\section{Classical Measured Stacking}
\label{sec:method-classical-stacking}

\subsection{Meta-Feature Construction}
\begin{itemize}
    \item Concatenate the measured probability outputs of all base VQCs.
    \item State the dimensionality of the resulting meta-feature vector.
    \item Explain whether original input features are included; use only the implemented final design.
\end{itemize}

\subsection{Out-of-Fold Training Procedure}
\begin{enumerate}
    \item Split the training data into stratified folds.
    \item Train a temporary ensemble on $K-1$ folds.
    \item Predict probabilities for the held-out fold.
    \item Repeat until every training sample has an out-of-fold meta-feature vector.
    \item Train the classical meta-learner on the complete out-of-fold matrix.
    \item Train the final base ensemble on the full training data.
    \item Generate final test meta-features with the final base ensemble.
\end{enumerate}

\subsection{Classical Meta-Learner}
\begin{itemize}
    \item State the final selected model, such as SVC or the implemented neural meta-learner.
    \item Include relevant hyperparameters and class weighting.
    \item Explain why this constitutes a strong measured baseline.
\end{itemize}

\section{Quantum Re-Encoding Stacking}
\label{sec:method-reencoding}

\subsection{Measured Inputs}
\begin{itemize}
    \item Use the same base-output information available to classical stacking.
    \item Explain scaling, padding or feature selection applied before quantum encoding.
\end{itemize}

\subsection{Meta-VQC}
\begin{itemize}
    \item Number of qubits and measured wires.
    \item Feature encoding of the base probability vector.
    \item Trainable and entangling layers.
    \item Loss, optimiser and training protocol.
\end{itemize}

\subsection{Purpose of the Condition}
\begin{itemize}
    \item Isolate whether a quantum meta-model adds value when the information supplied to it is already classical.
    \item Separate the effect of quantum meta-model capacity from the effect of preserving base states.
\end{itemize}

\section{Measurement-Free Quantum Stacking}
\label{sec:method-coherent}

\subsection{Architectural Principle}
\begin{itemize}
    \item Each trained VQC becomes an unmeasured subcircuit within one larger quantum circuit.
    \item Base learners occupy disjoint registers and process their assigned feature subsets.
    \item No base-level probability measurement is performed.
    \item A quantum meta-circuit acts across selected wires or registers.
    \item Measurement occurs only at the final classifier output.
\end{itemize}

\subsection{Base-Circuit Composition}
\begin{itemize}
    \item Reconstruct each selected ansatz inside the combined QNode.
    \item Load trained base parameters without copying the measured QLayer itself.
    \item Define the global wire mapping from local learner wires to the combined register.
    \item Preserve learner-specific feature subsets and circuit templates.
\end{itemize}

\subsection{Frozen Base Parameters}
\begin{itemize}
    \item Train each base learner before constructing the coherent circuit.
    \item Freeze the base parameters during initial meta-training.
    \item Train only meta-level parameters to reduce optimisation difficulty.
    \item State whether end-to-end fine-tuning was attempted; omit it if not executed.
\end{itemize}

\subsection{Meta-Wire Selection}
\begin{itemize}
    \item Main design: first or designated main wire from each base learner.
    \item Alternative design: all base wires, if evaluated.
    \item Explain the information--resource trade-off.
\end{itemize}

\subsection{Cross-Learner Meta-Circuit}
\begin{itemize}
    \item Meta-level RX/RZ rotations.
    \item CZ or RZZ connections between learner registers.
    \item Number of meta-layers and connection topology.
    \item How cross-register entanglement differs from independent measured aggregation.
\end{itemize}

\subsection{Final Measurement and Class Mapping}
\begin{itemize}
    \item Final measured wires.
    \item Basis-state to class-probability mapping.
    \item Treatment of unused basis states.
    \item Loss function used for final multiclass training.
\end{itemize}

\subsection{Scaling Analysis}
\begin{itemize}
    \item Total width approximately equals the sum of base-register widths.
    \item Base-circuit depth remains parallel conceptually, but simulator and transpiled depth depend on scheduling and topology.
    \item Meta-depth grows with meta-layer count and connectivity.
    \item Two-qubit gate count grows with base entanglement plus cross-learner meta-entanglement.
    \item Statevector memory scales exponentially with total coherent qubit count.
\end{itemize}

\researchnotes
\begin{itemize}
    \method{\item Staged optimisation of pretrained and frozen components is consistent with the broader hybrid variational approach and reduces the initial coherent optimisation burden~\cite{resch_quantum_2019,du_quantum_2025}.}
    \challenge{\item Cross-register two-qubit gates introduce the mechanism of interest but also increase depth, topology overhead and hardware noise exposure~\cite{resch_quantum_2019}.}
    \comparison{\item Multi-chip VQCs deliberately avoid inter-register entanglement for trainability, whereas this dissertation introduces a limited meta-circuit specifically to test coherent aggregation~\cite{park_addressing_2025}.}
\end{itemize}

\section{Classical Reference Models}
\label{sec:method-classical-baselines}

\sectionbrief
\begin{itemize}
    \item SVC configuration.
    \item Random Forest configuration.
    \item Preprocessing and data splits shared with quantum models.
    \item Explain that these provide practical context rather than perfectly capacity-matched quantum-advantage tests.
\end{itemize}

\section{Training Protocol}
\label{sec:method-training}

\subsection{Optimisation}
\begin{itemize}
    \item Adam or final selected optimiser.
    \item Learning rate.
    \item Negative log-likelihood or final selected loss.
    \item Batch size.
    \item Maximum epochs.
\end{itemize}

\subsection{Early Stopping}
\begin{itemize}
    \item Warm-up period.
    \item Patience.
    \item Minimum loss improvement.
    \item Restoration of the best model state.
\end{itemize}

\subsection{Repeated Training and Randomness}
\begin{itemize}
    \item Data-split seeds.
    \item Feature-subset and bagging seeds.
    \item Parameter initialisation.
    \item Number of repeated runs used in final comparisons.
    \item Explicitly acknowledge configurations with fewer repeats due to computational cost.
\end{itemize}

\subsection{Execution and Failure Handling}
\begin{itemize}
    \item Multiprocessing and HPC worker allocation.
    \item Memory and qubit-count limits.
    \item Failed-run logging.
    \item Rules for retrying, excluding or reporting failed configurations.
\end{itemize}

\section{Evaluation Framework}
\label{sec:method-evaluation}

\subsection{Predictive Metrics}
\begin{itemize}
    \item Accuracy.
    \item Macro-F1 as the primary class-balanced metric.
    \item Weighted-F1.
    \item Per-class precision and recall.
    \item Confusion matrices for selected comparisons.
\end{itemize}

\subsection{Optimisation Metrics}
\begin{itemize}
    \item Best and final training loss.
    \item Best and final validation loss.
    \item Epochs completed and early-stopping epoch.
    \item Convergence curves for representative configurations.
    \item Run failures or unstable training.
\end{itemize}

\subsection{Computational Metrics}
\begin{itemize}
    \item Base-training time.
    \item Out-of-fold stacking overhead.
    \item Meta-training time.
    \item End-to-end training time.
    \item Peak memory where available.
\end{itemize}

\subsection{Quantum Resource Metrics}
\begin{itemize}
    \item Qubit count.
    \item Circuit depth.
    \item Trainable parameter count.
    \item Total gate count.
    \item Two-qubit gate count.
    \item Transpiled depth and SWAP overhead, if transpilation is performed.
\end{itemize}

\researchnotes
\begin{itemize}
    \benchmark{\item Accuracy alone can obscure minority-class failure, supporting macro-F1 and per-class recall~\cite{demir_modified_2018,vasan_exploring_2025}.}
    \benchmark{\item Hardware feasibility depends on width, depth, two-qubit gate count, topology, gate fidelity and coherence rather than qubit count alone~\cite{resch_quantum_2019}.}
    \benchmark{\item Previous VQC ensemble studies evaluate performance and training cost jointly and use statistical tests where repeated observations permit them~\cite{radojicic_reducing_2026}.}
\end{itemize}

\section{Noise and Hardware Evaluation}
\label{sec:method-noise-hardware}

\sectionbrief
\begin{itemize}
    \item Keep this section only if noisy simulation or hardware experiments are completed.
    \item Specify the noise model or exact backend and calibration date.
    \item Specify shot count, transpilation level and error-mitigation method.
    \item Use a reduced set of configurations selected before hardware testing.
    \item Compare degradation relative to noiseless simulation.
    \item Do not retain generic IBM hardware claims if the final work remains simulator-only.
\end{itemize}

\researchnotes
\begin{itemize}
    \validation{\item Ensemble QNNs and AdaBoost.Q have been tested on superconducting hardware, showing that some quantum ensemble methods can move beyond ideal simulation~\cite{incudini_resource_2023,song_quantum_2024}.}
    \challenge{\item Hardware deployment requires topology-aware transpilation, qubit placement and error mitigation, and added two-qubit gates may remove a noiseless coherent advantage~\cite{rathi_enhanced_2025}.}
\end{itemize}

\section{Reproducibility and Result Management}
\label{sec:method-reproducibility}

\sectionbrief
\begin{itemize}
    \item Repository commit used for final experiments.
    \item Environment and dependency versions.
    \item Dataset download and preprocessing scripts.
    \item Configuration templates and job generator.
    \item Stable job IDs and skipped completed jobs.
    \item JSONL schema and reduced logging design.
    \item Aggregated parquet files and analysis dashboard.
    \item Random seeds and commands needed to reproduce the final tables.
\end{itemize}

% =============================================================================
% CHAPTER 4: EVALUATION AND RESULTS
% =============================================================================
\chapter{Evaluation and Results}
\label{chap:results}

\section{Experimental Overview}
\label{sec:results-overview}

\resultplaceholder
\begin{itemize}
    \item Final datasets and their roles.
    \item Final model configurations and aggregation conditions.
    \item Number of runs, seeds and completed/failed jobs.
    \item Hardware and simulator resources used.
    \item Compact table of the final experimental matrix.
\end{itemize}

\section{VQC Architecture Exploration}
\label{sec:results-vqc-search}

\subsection{Search Coverage}
\begin{itemize}
    \item Number of generated and completed VQC jobs.
    \item Distribution across datasets, qubit counts and architecture groups.
    \item Exclusions caused by failure, memory limits or invalid configurations.
\end{itemize}

\subsection{Overall Architecture Trends}
\begin{itemize}
    \item Effects of qubit count.
    \item Effects of feature density and features per qubit.
    \item Effects of re-uploading and trainable-layer count.
    \item Effects of entangler and entanglement placement.
    \item Effects of measurement mode.
    \item Report interactions only when supported by the analysis.
\end{itemize}

\subsection{Dataset-Dependent Behaviour}
\begin{itemize}
    \item Which architectural choices generalise across datasets.
    \item Which choices are strongly dataset-specific.
    \item Relationship with class count, feature count and dataset difficulty.
\end{itemize}

\subsection{Selected Base Configurations}
\begin{itemize}
    \item Final shortlist.
    \item Average percentile or ranking.
    \item Resource cost.
    \item Architectural diversity.
    \item Reason each design was retained for ensemble experiments.
\end{itemize}

\section{Baseline Classification Performance}
\label{sec:results-baselines}

\subsection{Classical Models}
\begin{itemize}
    \item SVC and Random Forest accuracy, macro-F1 and weighted-F1.
    \item Identify datasets where classical reference performance is already near ceiling.
\end{itemize}

\subsection{Individual VQCs}
\begin{itemize}
    \item Performance distribution of selected VQCs.
    \item Variance across architecture and seed.
    \item Resource cost and training time.
\end{itemize}

\subsection{Direct QuantumECOC}
\begin{itemize}
    \item Compare with individual VQCs.
    \item Evaluate the effect of ensemble size and base diversity.
    \item Report class-level gains or failures.
    \item Answer Research Question 1.
\end{itemize}

\section{Classical Measured Stacking Results}
\label{sec:results-classical-stacking}

\subsection{Predictive Performance}
\begin{itemize}
    \item Direct QuantumECOC versus classical StackedECOC.
    \item Accuracy and macro-F1 differences by dataset.
    \item Mean/median difference across datasets and seeds.
\end{itemize}

\subsection{Per-Class Behaviour}
\begin{itemize}
    \item Minority or difficult classes assisted by stacking.
    \item Confusion matrices for representative improvements and regressions.
\end{itemize}

\subsection{Computational Cost}
\begin{itemize}
    \item Cost of $K$ temporary fold ensembles.
    \item Final full-data retraining cost.
    \item Total stacking overhead relative to direct aggregation.
\end{itemize}

\subsection{Research Question 2}
\begin{itemize}
    \item Give a direct evidence-based answer after the results are presented.
\end{itemize}

\section{Quantum Re-Encoding Results}
\label{sec:results-reencoding}

\subsection{Predictive Performance}
\begin{itemize}
    \item Quantum re-encoding versus the classical meta-learner using the same measured base outputs.
    \item Dataset-level and aggregate comparisons.
\end{itemize}

\subsection{Convergence and Stability}
\begin{itemize}
    \item Loss curves.
    \item Sensitivity to initialisation.
    \item Early-stopping behaviour.
    \item Failed or collapsed runs.
\end{itemize}

\subsection{Research Question 3}
\begin{itemize}
    \item State whether making only the meta-model quantum provided a consistent benefit.
\end{itemize}

\section{Measurement-Free Stacking Results}
\label{sec:results-coherent}

\subsection{Implementation Feasibility}
\begin{itemize}
    \item Largest successfully simulated coherent circuit.
    \item Maximum learner and qubit counts.
    \item Memory, runtime and differentiation constraints.
    \item Configurations that could not be trained and the observed failure mode.
    \item Answer the feasibility component of Research Question 4.
\end{itemize}

\subsection{Predictive Performance}
\begin{itemize}
    \item Coherent stacking versus direct QuantumECOC.
    \item Coherent stacking versus classical stacking.
    \item Coherent stacking versus quantum re-encoding.
    \item Accuracy, macro-F1 and per-class behaviour.
\end{itemize}

\subsection{Convergence and Stability}
\begin{itemize}
    \item Training and validation curves.
    \item Effect of frozen base parameters.
    \item Meta-gradient and optimiser behaviour if analysed.
    \item Variability across seeds.
\end{itemize}

\subsection{Effect of Ensemble Size}
\begin{itemize}
    \item Predictive effect of increasing learner count.
    \item Qubit, depth, gate and simulation-cost growth.
    \item Point at which added learners stop being practical or beneficial.
\end{itemize}

\subsection{Effect of Meta-Circuit Design}
\begin{itemize}
    \item Main wires versus all wires, if evaluated.
    \item Meta-layer count.
    \item CZ versus RZZ.
    \item Entanglement topology.
\end{itemize}

\subsection{Dataset-Dependent Effects}
\begin{itemize}
    \item Association with class count, base weakness, imbalance or feature complexity.
    \item Distinguish robust patterns from isolated best scores.
\end{itemize}

\subsection{Research Questions 4 and 5}
\begin{itemize}
    \item Give direct answers on feasibility and predictive benefit.
\end{itemize}

\section{Resource and Scalability Analysis}
\label{sec:results-resources}

\subsection{Circuit Width}
\begin{itemize}
    \item Formula and measured values by ensemble size.
\end{itemize}

\subsection{Depth and Gate Counts}
\begin{itemize}
    \item Logical circuit depth.
    \item Total gates.
    \item Two-qubit gates.
    \item Transpiled values where available.
\end{itemize}

\subsection{Training Time and Memory}
\begin{itemize}
    \item Runtime scaling by qubit and learner count.
    \item Base training, fold training and coherent meta-training.
    \item Peak memory or practical cluster limits.
\end{itemize}

\subsection{Performance--Cost Trade-Off}
\begin{itemize}
    \item Performance gain per additional qubit, two-qubit gate or training hour.
    \item Identify configurations that are dominated by cheaper alternatives.
    \item Answer Research Question 6.
\end{itemize}

\section{Noise or Hardware Results}
\label{sec:results-noise}

\resultplaceholder
\begin{itemize}
    \item Retain only if the corresponding methodology was executed.
    \item Noiseless versus noisy simulation.
    \item Simulator versus hardware.
    \item Effect of transpilation and mitigation.
    \item Whether any coherent advantage survives realistic noise.
\end{itemize}

\section{Statistical and Practical Significance}
\label{sec:results-significance}

\sectionbrief
\begin{itemize}
    \item Confidence intervals or standard deviations across repeats.
    \item Paired comparisons using shared splits/seeds.
    \item Friedman/Wilcoxon or another justified test when enough observations exist.
    \item Effect sizes and consistency, not only $p$-values.
    \item Distinguish numerical best scores from practically meaningful improvements.
\end{itemize}

\section{Summary by Research Question}
\label{sec:results-rq-summary}

\begin{table}[ht]
    \centering
    \caption{Summary of evidence and findings by research question.}
    \label{tab:rq-summary}
    \begin{tabular}{p{0.08\textwidth}p{0.42\textwidth}p{0.42\textwidth}}
        \toprule
        RQ & Evidence & Finding \\
        \midrule
        RQ1 & Individual VQC and direct QuantumECOC results & \textit{To complete} \\
        RQ2 & Direct aggregation versus classical stacking & \textit{To complete} \\
        RQ3 & Classical meta-learner versus quantum re-encoding & \textit{To complete} \\
        RQ4 & Coherent implementation and scaling observations & \textit{To complete} \\
        RQ5 & Coherent versus measured predictive results & \textit{To complete} \\
        RQ6 & Width, depth, gates, time and memory & \textit{To complete} \\
        RQ7 & Dataset and ensemble characteristic analysis & \textit{To complete} \\
        \bottomrule
    \end{tabular}
\end{table}

% =============================================================================
% CHAPTER 5: DISCUSSION
% =============================================================================
\chapter{Discussion}
\label{chap:discussion}

\section{Interpretation of the Main Findings}
\label{sec:discussion-main}

\sectionbrief
\begin{itemize}
    \item Explain the overall pattern across direct aggregation, classical stacking, quantum re-encoding and coherent stacking.
    \item Identify whether performance was mainly determined by base-learner quality, learned aggregation, quantum meta-model capacity or state preservation.
    \item Interpret unsuccessful or unstable coherent experiments as evidence about feasibility rather than omitting them.
\end{itemize}

\section{Is Measurement an Information Bottleneck?}
\label{sec:discussion-measurement}

\sectionbrief
\begin{itemize}
    \item Return directly to the central hypothesis.
    \item Separate three possibilities: coherent improvement, no meaningful difference or worse coherent performance.
    \item Explain that poor coherent performance does not prove that the preserved state contains no useful information; optimisation, circuit design and resource limits may prevent its extraction.
    \item Avoid equating simulator state access with a practical hardware advantage.
\end{itemize}

\section{When Does Stacking Help?}
\label{sec:discussion-when}

\sectionbrief
\begin{itemize}
    \item Relate stacking gains to base-learner weakness.
    \item Examine class count, imbalance, feature dimensionality and dataset difficulty.
    \item Discuss whether simple datasets leave little room for a meta-learner.
    \item Discuss whether harder datasets benefit from learned aggregation but incur much higher training cost.
    \item Answer Research Question 7.
\end{itemize}

\section{Quantum Meta-Learning versus Coherent Aggregation}
\label{sec:discussion-meta-vs-coherent}

\sectionbrief
\begin{itemize}
    \item Contrast placing a quantum model after measurement with maintaining a continuous quantum computation across base and meta stages.
    \item Use the re-encoding condition to determine whether the meta-model type or the information boundary is more influential.
    \item Discuss capacity matching and optimisation fairness between classical and quantum stackers.
\end{itemize}

\section{Performance--Cost Trade-Off}
\label{sec:discussion-cost}

\sectionbrief
\begin{itemize}
    \item Evaluate whether any coherent performance gain justifies the extra qubits, depth, gates, memory and runtime.
    \item Discuss whether a smaller measured model is practically preferable even when coherent stacking achieves a slightly higher score.
    \item Identify the most promising configuration under a balanced performance--resource criterion.
\end{itemize}

\section{Relationship to Existing Literature}
\label{sec:discussion-literature}

\sectionbrief
\begin{itemize}
    \item Compare the observed ensemble gains with measured quantum bagging and boosting studies.
    \item Compare the classical stacking findings with measured quantum stacking literature.
    \item Compare the re-encoding findings with quantum meta-learning work.
    \item Compare coherent feasibility with theoretical or indexed coherent ensemble proposals.
    \item Explain which findings confirm, extend or conflict with earlier studies.
\end{itemize}

\section{Limitations}
\label{sec:discussion-limitations}

\sectionbrief
\begin{itemize}
    \item Small or moderate tabular datasets.
    \item Limited repeated runs for expensive configurations.
    \item Restricted hyperparameter search for the coherent condition.
    \item Exponential statevector simulation cost and resulting learner-count ceiling.
    \item Simulator-only conclusions if no hardware experiments are completed.
    \item Unequal representational capacity across classical and quantum meta-learners.
    \item Dependence on selected ansatz families, feature mapping and output measurement.
    \item Frozen-base training may prevent globally optimal end-to-end solutions.
    \item Accuracy and macro-F1 provide indirect rather than information-theoretic evidence about measurement loss.
    \item Internal student projects should be used as design guidance and clearly distinguished from peer-reviewed evidence~\cite{radojicic_reducing_2026,pinto_composing_2026}.
\end{itemize}

\section{Threats to Validity}
\label{sec:discussion-validity}

\subsection{Internal Validity}
\begin{itemize}
    \item Data leakage, preprocessing differences, random initialisation, feature selection and unequal optimisation budgets.
\end{itemize}

\subsection{Construct Validity}
\begin{itemize}
    \item Whether predictive metrics adequately represent preserved quantum information or information efficiency.
\end{itemize}

\subsection{External Validity}
\begin{itemize}
    \item Generalisation from small tabular datasets and simulation to larger datasets and physical devices.
\end{itemize}

\subsection{Conclusion Validity}
\begin{itemize}
    \item Number of repeats, failed runs, statistical power and sensitivity to selected datasets.
\end{itemize}

% =============================================================================
% CHAPTER 6: CONCLUSION
% =============================================================================
\chapter{Conclusion and Future Work}
\label{chap:conclusion}

\section{Dissertation Summary}
\label{sec:conclusion-summary}

\sectionbrief
\begin{itemize}
    \item Restate the problem and motivation.
    \item Summarise the progression from VQC selection and ECOC construction to the three stacking conditions.
    \item Summarise the final experimental scope.
\end{itemize}

\section{Answers to the Research Questions}
\label{sec:conclusion-rqs}

\sectionbrief
\begin{itemize}
    \item Answer each research question explicitly in order.
    \item Use concise evidence-based statements rather than repeating the whole results chapter.
    \item State uncertainty and conditional conclusions where the experiments are limited.
\end{itemize}

\section{Contributions}
\label{sec:conclusion-contributions}

\sectionbrief
\begin{itemize}
    \item Restate only completed implementation contributions.
    \item State empirical findings separately from software or methodological contributions.
    \item Include informative negative results, practical scaling limits and failed approaches where they advance understanding.
\end{itemize}

\section{Overall Conclusion}
\label{sec:conclusion-overall}

\sectionbrief
\begin{itemize}
    \item Give the final judgement on whether measurement-free stacking was feasible.
    \item State whether it offered a consistent predictive advantage.
    \item State whether any gain justified its additional resource and optimisation cost.
    \item Avoid broader quantum-advantage claims beyond the evidence.
\end{itemize}

\section{Future Work}
\label{sec:conclusion-future}

\sectionbrief
\begin{itemize}
    \item End-to-end fine-tuning after frozen-base meta-training.
    \item Alternative cross-register meta-circuits and learned connectivity.
    \item More direct diagnostics of phase, amplitude or mutual-information preservation.
    \item Larger coherent ensembles using tensor-network, distributed or approximate simulation.
    \item Hardware-aware training, topology-constrained meta-circuits and error mitigation.
    \item Real IBM hardware evaluation of a reduced coherent configuration.
    \item Stronger repeated-seed and statistical evaluation.
    \item Capacity-matched monolithic VQC and classical neural meta-learner comparisons.
    \item Automated selection based on prediction diversity rather than architecture labels alone.
    \item Release or formal evaluation of the VQC architecture-search pipeline as a secondary research tool.
\end{itemize}

% =============================================================================
% APPENDICES
% =============================================================================
\appendix

\chapter{Dataset Details}
\label{app:datasets}
\begin{itemize}
    \item Data sources and licences.
    \item Original and cleaned dimensions.
    \item Feature definitions and class distributions.
    \item Dataset-specific preprocessing.
\end{itemize}

\chapter{VQC and Ansatz Configurations}
\label{app:ansatze}
\begin{itemize}
    \item Layer specifications.
    \item Final template JSON definitions.
    \item Parameter and gate counts.
    \item Representative circuit diagrams.
\end{itemize}

\chapter{Experimental Search Spaces}
\label{app:search-spaces}
\begin{itemize}
    \item VQC search grid.
    \item Ensemble composition grid.
    \item Meta-circuit grid.
    \item Final reduced stress-test configurations.
\end{itemize}

\chapter{Additional Results}
\label{app:additional-results}
\begin{itemize}
    \item Full per-dataset tables.
    \item Classification reports.
    \item Confusion matrices.
    \item Loss curves.
    \item Failed-run summaries.
\end{itemize}

\chapter{Circuit and Resource Details}
\label{app:circuits}
\begin{itemize}
    \item Individual VQC circuit.
    \item Direct QuantumECOC information flow.
    \item Re-encoding meta-VQC.
    \item Measurement-free coherent circuit.
    \item Logical and transpiled resource tables.
\end{itemize}

\chapter{Computational Environment and Reproducibility}
\label{app:reproducibility}
\begin{itemize}
    \item Software versions.
    \item HPC resources and Slurm scripts.
    \item Commands for job generation, execution and aggregation.
    \item Result schemas and dashboard instructions.
\end{itemize}

% =============================================================================
% ACKNOWLEDGEMENTS AND BIBLIOGRAPHY
% =============================================================================
\chapter*{Acknowledgements}
% Write after the main dissertation is complete.
% Suggested people and resources to consider: Dr Simon Caton, programme staff,
% UCD SONIC support, family and anyone who directly supported the project.

\printbibliography

\end{document}
